% Generated draft: requires mathematical and source review.
\documentclass[11pt]{article}
\usepackage[margin=25mm]{geometry}
\usepackage{fontspec}
\setmainfont{Latin Modern Roman}
\usepackage{amsmath,amssymb,mathtools,unicode-math}
\setmathfont{Latin Modern Math}
\usepackage{enumitem,needspace,xurl,hyperref}
\setlength{\parindent}{0pt}
\setlength{\parskip}{5pt}
\setlength{\emergencystretch}{3em}
\newcommand{\N}{\mathbb{N}}
\newcommand{\Z}{\mathbb{Z}}
\newcommand{\Q}{\mathbb{Q}}
\newcommand{\R}{\mathbb{R}}
\newcommand{\C}{\mathbb{C}}
\newcommand{\F}{\mathbb{F}}
\newcommand{\A}{\mathbb{A}}
\newcommand{\PP}{\mathbb{P}}
\newcommand{\E}{\mathbb{E}}
\newcommand{\eps}{\varepsilon}
\newcommand{\dd}{\,\mathrm{d}}
\begin{document}
\section*{UnsolvedMath ID: 30004892}
\textbf{Source problem code:} catalogue entry
\subsection{Definitions and mathematical statement}
The supplied record concerns a conjectural hypercube-cluster formula, identified by its title as relating to Kazhdan--Lusztig polynomials. It does not supply the formula itself or complete definitions of its hypotheses. Consequently, the following is a schematic statement of the recorded target, requiring reconstruction from the cited report before it can be a self-contained mathematical conjecture.

For an ambient partially ordered set \((P,\leq)\), interval notation means
\[
[x,y]=\{z\in P:x\leq z\leq y\}.
\]
The intended ambient ordered set, graph structure, and restrictions on the endpoints \(x\leq y\) are not specified in the supplied statement. The dated assessment describes the unresolved assertion as being in type \(A\); this does not by itself specify these missing data.

With the intended ambient data and endpoints fixed, the recorded conjecture asks whether, for every \(c'\in[x,y]\), the proposed formula holds whenever both of the following conditions are satisfied:
\begin{enumerate}
\item The subinterval \([x,c']\) is diamond complete, in the sense attributed in the excerpt to Patimo, reference [2].
\item For every vertex \(v\in[x,c']\), the edges leaving \([x,c']\) from \(v\) span a hypercube cluster.
\end{enumerate}
Equivalently, its quantifier structure is schematically
\[
\forall c'\in[x,y],\quad
\left(
\begin{gathered}
[x,c']\text{ is diamond complete},\\
\text{and for every }v\in[x,c'],\\
\text{the edges leaving }[x,c']\text{ from }v\\
\text{span a hypercube cluster}
\end{gathered}
\right)
\Longrightarrow
\left(\text{the proposed hypercube-cluster formula holds for }c'\right).
\]
Here the technical predicates and the conclusion retain their source-dependent meanings; they have not been defined by this display. The excerpt informally connects the hypercube-cluster condition with existence of least common multiples, but gives neither the relevant order nor the collections for which these must exist. It also omits the graph and its edge convention, and whether diamond completeness is relative to the interval or a larger graph.

The extracted geometric background mentions a quotient of a punctured object by an attractive torus action, and an action fixing \(\mathrm{GL}_n/B\subset\mathrm{GL}_{n+1}/B\) while acting nontrivially on the remainder, producing an open/closed decomposition. The notation for the punctured object and its quotient is corrupted, and the decomposition and its relationship to the conjectural formula are not supplied. These fragments do not determine additional precise hypotheses.
\subsection{Short English statement}
The recorded conjecture asks whether the proposed hypercube-cluster formula holds for every eligible \(c'\in[x,y]\): the interval \([x,c']\) must be diamond complete, and the edges leaving it from each of its vertices must span a hypercube cluster. The formula and essential technical definitions are missing, so this remains a schematic, statement-only draft requiring human mathematical review.
\subsection{Sources}
\begin{itemize}
\item Local registry: \texttt{data/unsolvedmath/problems.json}, with \texttt{catalog/problems.json} as the additional-record registry; supplied permanent UnsolvedMath ID \(30004892\), source problem code \texttt{catalogue entry}, title \emph{Diamond-Complete Hypercube Formulas for Kazhdan--Lusztig Polynomials}. These registry references identify the supplied record; no independent file inspection is asserted.
\item Supplied report citation: \emph{Enveloping Algebras and Geometric Representation Theory} (2022), Oberwolfach Reports, pp.\ 2827--2891, DOI \texttt{10.4171/owr/2021/52}. \url{https://doi.org/10.4171/owr/2021/52}. The record gives proposed year 2021 but names no proposer. The excerpt attributes the notion of diamond completeness to Patimo, reference [2], without supplying that reference's bibliographic details.
\item Supplied literature source, labeled ``BBDVW lower-interval theorem'': \url{https://doi.org/10.1093/imrn/rnag022}.
\item Supplied literature source, labeled ``Zannoni, standard decompositions'': \url{https://arxiv.org/abs/2605.13304}.
\item Supplied literature source, labeled ``Gurevich--Wang, hypercube decomposition'': \url{https://link.springer.com/article/10.1007/s00029-024-00972-0}.
\item Literature assessment dated exactly \texttt{2026-08-22}: the supplied curation labels the problem partially solved and states that the formula is proved for lower intervals and several standard or stronger subclasses, while the universal type-\(A\) assertion remains unresolved. This is supplied dated curation, not an independently verified literature assessment; its status and cited claims require independent verification.
\item The registry labels the statement \texttt{corrected\_verified} and reports that the official report was checked directly and that definitions were restored. Those are supplied provenance claims. The present evidence does not include the missing formula or complete definitions, and no source visit or independent verification is asserted here.
\end{itemize}
\end{document}
